\section{Methods}
\label{sec:methods}

We consider the standard linear regression model
\begin{equation}
\label{eq:linear_model}
\by = \bX \bbeta + \beps, \qquad \beps \sim \Normal(\mathbf{0}, \sigma^2 \mathbf{I}_n),
\end{equation}
where $\bX \in \R^{n \times p}$ is the design matrix, $\bbeta \in \R^p$ is the coefficient vector, and $\sigma^2$ is the noise variance.
We assume $\bbeta$ is \emph{sparse}: only $s \ll p$ entries are nonzero.
The goal is to estimate $\bbeta$, predict on held-out data, and recover the support $S = \{j : \beta_j \neq 0\}$.

Below we describe the six methods included in our benchmark, grouped into classical penalized estimators and Bayesian approaches.

% ------------------------------------------------------------------
\subsection{Classical Penalized Methods}
% ------------------------------------------------------------------

\subsubsection{Ordinary Least Squares (OLS)}

OLS minimizes the unpenalized squared loss:
\begin{equation}
\hat{\bbeta}_{\text{OLS}} = \arg\min_{\bbeta} \frac{1}{2n} \|\by - \bX \bbeta\|_2^2 = (\bX^\top \bX)^{-1} \bX^\top \by.
\end{equation}
OLS serves as an unregularized baseline.
It is consistent when $n > p$ but does not perform variable selection and is undefined when $p > n$.

\subsubsection{Ridge Regression}

Ridge regression \citep{hoerl1970ridge} adds an $\ell_2$ penalty:
\begin{equation}
\hat{\bbeta}_{\text{Ridge}} = \arg\min_{\bbeta} \frac{1}{2n} \|\by - \bX \bbeta\|_2^2 + \lambda \|\bbeta\|_2^2.
\end{equation}
The penalty shrinks coefficients toward zero but does not set any exactly to zero.
From a Bayesian perspective, this corresponds to a Gaussian prior $\beta_j \sim \Normal(0, 1/(2\lambda))$.
The regularization parameter $\lambda$ is selected by leave-one-out cross-validation over a grid of 50 logarithmically spaced values in $[10^{-4}, 10^{4}]$.

\subsubsection{Lasso}

The Lasso \citep{tibshirani1996regression} uses an $\ell_1$ penalty:
\begin{equation}
\hat{\bbeta}_{\text{Lasso}} = \arg\min_{\bbeta} \frac{1}{2n} \|\by - \bX \bbeta\|_2^2 + \lambda \|\bbeta\|_1.
\end{equation}
The $\ell_1$ penalty induces sparsity, setting small coefficients exactly to zero.
Its Bayesian interpretation is a Laplace (double exponential) prior on each $\beta_j$.
We use 5-fold cross-validation over 100 candidate $\lambda$ values.

\subsubsection{Elastic Net}

The Elastic Net \citep{zou2005regularization} combines both penalties:
\begin{equation}
\hat{\bbeta}_{\text{EN}} = \arg\min_{\bbeta} \frac{1}{2n} \|\by - \bX \bbeta\|_2^2 + \lambda \left[\alpha \|\bbeta\|_1 + (1-\alpha) \|\bbeta\|_2^2\right],
\end{equation}
where $\alpha \in [0,1]$ controls the mix between $\ell_1$ and $\ell_2$ penalties.
This encourages sparsity while handling correlated features better than pure Lasso through the grouping effect.
We select both $\lambda$ and $\alpha$ via 5-fold cross-validation, testing $\alpha \in \{0.1, 0.5, 0.7, 0.9, 0.95, 0.99, 1.0\}$.

% ------------------------------------------------------------------
\subsection{Bayesian Sparse Priors}
% ------------------------------------------------------------------

Both Bayesian methods use the NUTS sampler \citep{hoffman2014no} as implemented in PyMC \citep{abril2023pymc}, with 2 chains, 1{,}000 warmup iterations, and 2{,}000 posterior draws per chain, with a target acceptance rate of 0.95.
Point estimates are posterior means; credible intervals are 95\% highest density intervals (HDI) computed via ArviZ \citep{kumar2019arviz}.

\subsubsection{Horseshoe Prior}

The Horseshoe prior \citep{carvalho2010horseshoe} uses a global-local shrinkage hierarchy:
\begin{align}
\beta_j \mid \lambda_j, \tau &\sim \Normal(0, \lambda_j^2 \tau^2), \quad j = 1, \ldots, p, \\
\lambda_j &\sim \HalfCauchy(0, 1), \label{eq:hs_local} \\
\tau &\sim \HalfCauchy(0, \tau_0), \label{eq:hs_global} \\
\sigma &\sim \HalfCauchy(0, 2),
\end{align}
where $\lambda_j$ is the local shrinkage parameter for coefficient $j$ and $\tau$ is the global shrinkage parameter.
The half-Cauchy prior on $\lambda_j$ has heavy tails that allow large signals to escape shrinkage while concentrating small signals near zero.
We set $\tau_0 = 1$ as the scale of the global shrinkage prior.

\subsubsection{Spike-and-Slab Prior}

We use a continuous relaxation of the Spike-and-Slab prior \citep{mitchell1988bayesian} to enable gradient-based NUTS sampling:
\begin{align}
\beta_j \mid z_j &\sim z_j \cdot \Normal(0, \sigma_{\text{slab}}^2) + (1 - z_j) \cdot \Normal(0, \sigma_{\text{spike}}^2), \\
z_j &\sim \text{Bernoulli}(\pi), \\
\pi &\sim \text{Beta}(1, 1/\pi_0), \\
\sigma &\sim \HalfCauchy(0, 2),
\end{align}
where $\sigma_{\text{slab}} = 5$ is the scale of the ``slab'' component (allowing large coefficients), $\sigma_{\text{spike}} = 0.01$ is the scale of the ``spike'' component (concentrating near zero), and $\pi_0 = 0.2$ is the prior expected inclusion probability.
The continuous mixture formulation via \texttt{NormalMixture} in PyMC enables efficient NUTS sampling without discrete sampling steps.
